(1) 面电荷密度 \(\sigma=Q/A\)。无限大导体板各自产生 \(\sigma/(2\epsilon_0)\) 场，两板内部叠加得 \(E=Q/(\epsilon_0 A)\)，外部抵消。\(V=Ed=Qd/(\epsilon_0A)\)、\(C=Q/V=\epsilon_0 A/d\)。

(2) 固定 Q 的场能 \(U=Q^2/(2C)=Q^2d/(2\epsilon_0 A)\)。以 d 增大为正坐标，场对极板的广义力 \(F_d=-\partial U/\partial d=-Q^2/(2\epsilon_0 A)\)，所以吸引压力
\[
p=\frac{|F_d|}{A}=\frac{\sigma^2}{2\epsilon_0}=\frac{\epsilon_0E^2}{2}.
\]

(3) 一块极板所受力来自另一块极板的电场，而不应把自身场算进去。另一块板产生的场大小 \(|E_{\rm other}|=\sigma/(2\epsilon_0)\)，所以该极板的力大小 \(Q|E_{\rm other}|=Q^2/(2\epsilon_0A)\)，与能量法一致。

(4) 外力准静态增大间距 \(d_1\to d_2\) 时对体系做功 \(W_{\rm ext}=U(d_2)-U(d_1)=Q^2(d_2-d_1)/(2\epsilon_0A)\)。若固定 V，板间电荷会随 d 变化，必须把电源的功加入能量预算，不能仅对 \(CV^2/2\) 取负导数。
